\subsection{Capability acquisition needs a negative-side audit as well as a positive delta}

A final adversarial-prior-art route surfaced a particularly close comparator for Self-\RAKL{}. Chen et al. argue that a benchmark delta cannot by itself distinguish a genuinely new operational capability from extra search, a changed verifier or adaptation to a fallible oracle, and instantiate a two-sided audit in which the prior verifier's representational range is bounded independently of the improved system \cite{chen2026twosided}. Their RNA case also shows how a result that appears strong under the optimized oracle can collapse under an outside predictor panel. This is directly consonant with Orion's evaluator-separation and fresh-assurance requirements, and it narrows any novelty claim based merely on ``held-out improvement.''

The comparison sharpens the Self-Orion receipt. A strong method-evolution claim should report separately: (i) what the incumbent could or could not do under a frozen pre-change evaluator and resource envelope; (ii) what the challenger does under an adjudicator outside the optimization loop; (iii) whether the difference survives matched-search/compute controls; and (iv) whether the hypothesized mechanism of improvement transfers under intervention. A positive development delta can satisfy none of these automatically. Where a decidable negative-side exclusion is available, it is stronger than a failed finite search and should be recorded as such. Where only an empirical floor is available, Orion must not promote it to impossibility.

This addition changes the novelty boundary rather than the architecture: protected assurance, evaluator identity, matched resources and mechanism transfer were already registered. The newly assimilated work provides a close independent precedent for why those coordinates must remain separate and for why a self-improvement claim needs evidence about the incumbent's prior capability boundary, not only the challenger's final score.
